\section*{Popular Summary}

Quantum circuits built from language structure have a special shape. A sentence, paragraph, dialogue, or document is not usually a random pattern of interactions. Words combine into phrases, phrases combine into larger units, and discourse relations connect pieces of text. In categorical and tensor-network approaches to language, this structure becomes a diagram. In quantum natural language processing it can become a quantum circuit or tensor network. The graph of that circuit matters for simulation.

There are two familiar reasons why a quantum circuit may be easy to simulate classically. The first is small graph width. If a circuit can be cut into pieces that communicate through small separators, tensor-network methods can contract it efficiently. The second is low magic. Stabilizer circuits, built from Clifford operations, stabilizer states, Pauli measurements, and related free components, are efficiently simulable. Non-Clifford operations introduce magic, and stabilizer-decomposition or quasiprobability simulations pay an overhead for that resource.

Width alone and magic alone can both be misleading. A width bound treats all tensors crossing a separator as dense objects, even when much of the network is stabilizer structure that can be represented symbolically. A global magic bound charges for all non-Clifford tensors at once, even when many of them live inside small modules whose effects are locally measured or marginalized before reaching high-level separators. For structured language circuits, where the location of a word or discourse update is part of the model, this distinction is important.

This paper introduces an explicit message-recurrence formulation of separator-local magic. The key quantity is no longer defined by saying that some good separator support set exists. Instead, the paper defines an explicit message recurrence. Each bag in a rooted tree decomposition has a local quasiprobability cost. Each child subtree is either active, meaning that its unresolved resource randomness is still carried by the outgoing message, or inactive, meaning that it has been locally compressed or marginalized into a bounded free or dense boundary message. The resulting burden \(\Lambda_{\mathrm{msg}}\) is the largest message cost produced by this recurrence.

This makes the theorem narrower but more defensible. The proved result applies to locally marginalizable resource-decomposition networks. These are structured circuits for which local non-Clifford choices can be absorbed, averaged, or compressed before they propagate through the whole graph. The theorem does not say that every QNLP or QDisCoCirc circuit has this property. Some circuits may deliberately place non-Clifford resources at high-level separators, and then the separator-local burden becomes large.

The paper also separates two notions of width. Ordinary tensor-network contraction depends on a dense width \(w_{\mathrm{TN}}\). The SLMS simulator depends on an effective width \(w_{\mathrm{eff}}\), the number of boundary degrees of freedom that must still be represented densely after stabilizer compression. A stabilizer graph state on a two-dimensional grid is a simple example: as a dense tensor network it has large grid treewidth, but as a stabilizer object it can be represented symbolically. The paper uses this kind of example to show that dense treewidth and stabilizer-compressed width can disagree. If \(w_{\mathrm{eff}}=w_{\mathrm{TN}}\) and the circuit already has small dense treewidth, ordinary tensor-network contraction may be better. The separator-local method is useful only when stabilizer compression or local resource marginalization gives a better effective parameter than dense contraction or global magic sampling.

The paper is deliberately conservative. It does not claim unconditional quantum advantage. It does not claim that quantum methods improve legal or other language applications. It also does not claim to replace modern tensor-network, ZX-calculus, stabilizer-decomposition, or stabilizer tensor-network simulators. Legal-text-derived graphs appear only as possible structured stress tests. The main purpose is to give a rigorous way to ask when structured quantum language circuits are classically easy, and which placements of non-stabilizer resource make them harder to simulate under this particular certificate.
